% Source: 01 - Mon  Sep 28 2026.pdf, page 7. Content remains in source order.
\sourcepage{01}{7}
\sourceheading{COMPLEXITY CLASSES AND THEORY OF INTRACTABILITY (NP-COMPLETENESS)}
\textbf{Objective:} CATEGORIZE computational problems according to the complexity of the algorithms for their solution.

\textbf{RECAP: Computational Problem}
\[
\Pi\subseteq\mathcal I\times\mathcal S\quad\text{(relation)},\qquad
(i,s)\in\Pi\quad[i\mathrel\Pi s]\quad\Longleftrightarrow\quad s\text{ is a solution to }i.
\]
Many-to-many!

Algorithm that solves $\Pi$: implements
\[
A_\Pi:\mathcal I\longrightarrow\mathcal S,\qquad
\forall i\in\mathcal I:\quad i\mathrel\Pi A_\Pi(i)\in\mathcal S.
\]
We are interested in upper/lower bounds (on algorithms) for $\Pi$:
\[
\underbrace{\exists A_\Pi\text{ solving }\Pi:\ T_A(n)=O(f(n))}_{\text{UPPER BOUND}},
\]
\[
\underbrace{\forall A_\Pi\text{ solving }\Pi:\ T_A(n)=\Omega(f(n))}_{\text{LOWER BOUND (HARDER!)}}.
\]
The order of $f(n)$ captures the difficulty of $\Pi$.
